\section{符号说明}
为避免符号表挤占正文并降低可读性，本文只列出后续推导反复使用的核心符号；仅出现一次的辅助量在公式后就地解释。任务、能源和评价指标分组列示，均采用模板默认正文字号。

\subsection{任务、预测与容量符号}
\begin{longtable}{p{0.20\textwidth}p{0.56\textwidth}p{0.14\textwidth}}
\caption{任务与预测核心符号}\label{tab:symbol-task}\\
\toprule 符号 & 含义 & 单位/类型\\ \midrule\endfirsthead
\toprule 符号 & 含义 & 单位/类型\\ \midrule\endhead
\bottomrule\endfoot
$i,r,t,k$ & 任务、区域、小时、任务类型索引 & 索引\\
$a_i,h_i,F_i$ & 到达时刻、持续小时数、有效截止时点 & h\\
$g_i,o_i,k_i$ & GPU 需求、来源区域、任务类型 & GPU/类别\\
$\ell_{o_ir},\ell_i^{\max}$ & 来源至执行区单向时延及任务上限 & ms\\
$\omega_{it}(s)$ & 开工时刻为 $s$ 时，任务与小时 $t$ 的重叠时长 & h\\
$\chi_{it}(s)$ & 任务在小时 $t$ 是否存在正重叠 & 二元参数\\
$\Omega_i$ & 任务 $i$ 的合法区域—开工候选集合 & 集合\\
$x_{irs}$ & 任务是否在区域 $r$、时刻 $s$ 开工 & 二元变量\\
$D_{rkt},\widehat D_{rkt}$ & 实际与预测的逐时到达 GPU 需求 & GPU\\
$x_{rk,t}$ & Huber 模型的滞后量与周期特征向量 & 向量\\
$G_r^{\max}$ & 区域可调度 GPU 容量 & GPU\\
$P_r^{IT,\max},P_r^{F,\max}$ & 区域 IT 与设施功率上限 & MW\\
$p_k^{GPU}$ & 任务类型单位 GPU 的平均 IT 功率 & MW/GPU\\
$\pi_r$ & 区域 PUE & 无量纲\\
$L_{rt}^{N},L_{rt}^{BAI}$ & 非 AI 负荷与问题三固定 AI 基准负荷 & MW\\
$L_{rt}^{AI,avg},L_{rt}^{F}$ & AI 小时平均负荷与设施负荷 & MW\\
\end{longtable}

\subsection{电网、新能源与储能符号}
\begin{longtable}{p{0.20\textwidth}p{0.56\textwidth}p{0.14\textwidth}}
\caption{能源与储能核心符号}\label{tab:symbol-energy}\\
\toprule 符号 & 含义 & 单位/类型\\ \midrule\endfirsthead
\toprule 符号 & 含义 & 单位/类型\\ \midrule\endhead
\bottomrule\endfoot
$R_{rt}^{av}$ & 可用新能源功率 & MW\\
$R_{rt}^{use},R_{rt}^{ch}$ & 新能源直接消纳与储能充电功率 & MW\\
$R_{rt}^{sell},R_{rt}^{curt}$ & 新能源外送与弃电功率 & MW\\
$G_{rt}^{buy},G_{rt}^{sell}$ & 电网购电与系统边界外送功率 & MW\\
$I_r^{\max},X_r^{grid}$ & 区域最大购电与最大外送功率 & MW\\
$C_{rt},D_{rt}$ & 储能充电与放电功率 & MW\\
$S_{rt},S_r^0$ & 时段末 SOC 与运行前初始 SOC & MWh\\
$S_r^{min},S_r^{max}$ & SOC 下限与上限 & MWh\\
$C_r^{\max},D_r^{\max}$ & 最大充电与放电功率 & MW\\
$\eta_r^c,\eta_r^d$ & 充电效率与放电效率 & 无量纲\\
$c_{rt}^{buy},c_{rt}^{sell}$ & 购电价格与售电价格 & 元/MWh\\
$\kappa_{rt}$ & 电网购电碳强度 & tCO$_2$/MWh\\
$N_{rt}$ & 净购电功率 $G_{rt}^{buy}-G_{rt}^{sell}$ & MW\\
\end{longtable}

\subsection{评价与搜索符号}
\begin{longtable}{p{0.20\textwidth}p{0.56\textwidth}p{0.14\textwidth}}
\caption{评价指标与搜索符号}\label{tab:symbol-metric}\\
\toprule 符号 & 含义 & 单位/类型\\ \midrule\endfirsthead
\toprule 符号 & 含义 & 单位/类型\\ \midrule\endhead
\bottomrule\endfoot
$C_{op}$ & 购电支出减售电收入后的运行成本 & 元\\
$E_{CO_2}$ & 电网购电产生的累计碳排放 & tCO$_2$\\
$U_{RE}$ & 新能源有效利用率 & 无量纲\\
$L_{net}$ & GPU-hour 加权网络时延 & ms\\
$Q_{service}$ & 实时、按期和等待指标形成的服务质量 & 无量纲\\
$P^{peak}$ & 六区域非负峰值净购电功率之和 & MW\\
$V$ & 相邻小时净购电绝对变化量之和 & MW\\
$N^{cycle}$ & 储能等效完整循环次数之和 & 次\\
$\Psi$ & 有限候选比较使用的标准化代理分数 & 无量纲\\
$\rho_C,\delta_P,\delta_S,\delta_R$ & 碳约束、峰谷价差、售电与新能源情景参数 & 无量纲\\
\end{longtable}

其中 $\chi_{it}$ 与 $\omega_{it}$ 分别服务于瞬时容量和小时电量，二者不可互换；$U_{RE}$、$L_{net}$ 和服务质量只有在分母有效时才进入评价。完整实现中的辅助变量、状态哈希与停止计数见程序及结构化运行证据。
